\documentclass[../../main.tex]{subfiles}
\begin{document}

\begin{center}
  {\fontsize{48pt}{58pt}\selectfont \textbf{1997}\normalsize\textbf{年度}}
\end{center}

\begin{table}[htb]
  \centering
  \caption{出題テーマと難易度}
  \begin{tabular}{|c|p{9cm}|c|}
    \hline
    \textbf{問題番号} & \textbf{テーマ} & \textbf{難易度} \\
    \hline
    第1問 & 正方形に含まれる正三角形の面積 & \\
    \hline
    第2問 & 二次式がすべての整数で正となる条件 & \\
    \hline
    第3問 & 空間の点の軌跡と内積の極限 & \\
    \hline
    第4問 & 正三角形内の光線の反射 & \\
    \hline
    第5問 & 不等式領域を回転してできる立体の体積 & \\
    \hline
    第6問 & 三次曲線と放物線の共通接線 & \\
    \hline
  \end{tabular}
\end{table}

\begin{table}[htb]
  \centering
  \caption{解答}
  \renewcommand{\arraystretch}{1.4}
  \begin{tabular}{|c|c|p{9cm}|}
    \hline
    \textbf{問題番号} & \textbf{小問} & \textbf{答え} \\
    \hline
    第1問 & (1) & $(a,b)$の範囲を図示（解答参照） \\
    \cline{2-3}
    & (2) & $(a,b)=(1,2-\sqrt{3}),
    (2-\sqrt{3},1),(\sqrt{3}-1,\sqrt{3}-1)$で
    最大値$2\sqrt{3}-3$ \\
    \hline
    第2問 & -- & $0<a<2n+1$ \\
    \hline
    第3問 & (1) & $\dfrac{\sqrt{2}}{2}\le r<1$ \\
    \cline{2-3}
    & (2) & $\displaystyle\lim_{r\to1-0}(1-r)^2
    \{M(r)-m(r)\}=\dfrac{1}{2}$ \\
    \hline
    第4問 & (1) & 頂点$B$に到達する \\
    \cline{2-3}
    & (2) & 頂点$C$に到達し，反射は$12k+3$回 \\
    \hline
    第5問 & -- & $V=\pi^2\left(\dfrac{1}{4}-a\right)$ \\
    \hline
    第6問 & (1) & $a>-\dfrac{1}{2}$かつ$a\ne0$ \\
    \cline{2-3}
    & (2) & $S=\dfrac{16}{3}(1+2a)^{\frac{3}{2}}$ \\
    \hline
  \end{tabular}
\end{table}

\paragraph{試験形態}
6問．試験時間は要確認．

\end{document}
